\begin{statementbox}
\textbf{\textcolor{CStatement}{C020｜椭圆切线距离与焦点距离乘积定理}}

\medskip
\textbf{\textcolor{CStatement}{条件}}\par
设椭圆
\[
E:\quad \frac{x^2}{a^2}+\frac{y^2}{b^2}=1,\qquad a>b>0,
\]
焦点为 $F_1(-c,0),F_2(c,0)$，其中 $c^2=a^2-b^2$。
任取椭圆上一点 $P(x_0,y_0)$，过 $P$ 作切线 $l$。
记
\[
r_1=|PF_1|,\qquad r_2=|PF_2|,\qquad
 d=\operatorname{dist}(O,l),\quad O=(0,0).
\]

\textbf{\textcolor{CStatement}{核心结论（任意切点均成立）}}
\[
\boxed{\sqrt{r_1r_2}\,d=ab}\qquad\Longleftrightarrow\qquad
\boxed{r_1r_2d^2=a^2b^2}.
\]
\textbf{一句话：}两条焦半径的\textbf{几何平均数}，乘以\textbf{椭圆中心到切线的垂距}，等于半长轴与半短轴的乘积。

\medskip
\textbf{\textcolor{CStatement}{最有用的等价形式}}
\[
\boxed{d=\frac{ab}{\sqrt{r_1r_2}}},\qquad
\boxed{r_1r_2=\frac{a^2b^2}{d^2}}.
\]
若令 $\theta=\angle F_1PF_2$，还有
\[
\boxed{d=a\cos\frac{\theta}{2}},\qquad
\boxed{r_1r_2\cos^2\frac{\theta}{2}=b^2}.
\]

\textbf{\textcolor{CStatement}{范围与特殊点}}
\[
b\le d\le a,\qquad b^2\le r_1r_2\le a^2.
\]
当 $P=(\pm a,0)$ 时，$d=a,\ r_1r_2=b^2$；
当 $P=(0,\pm b)$ 时，$d=b,\ r_1r_2=a^2$。
上述式子在两端点同样成立。

\textbf{\textcolor{CStatement}{使用边界}}\par
此处 $d$ 是\textbf{椭圆中心到该切线的距离}，不是切点到焦点的距离，也不是焦点到切线的距离。
对平移或旋转后的椭圆，应以其\textbf{中心}与\textbf{真实焦点}定义这些距离；公式不变。
\end{statementbox}
